% from ega2-1-fr.tex
\oldpage[II]{23}
étant, si $g\in S_m$, $x\in\mathfrak{p}_n$, on a
$f^rgx\in\mathfrak{p}_{m+n+rd}$ sauf pour un nombre fini de valeurs de $r$,
donc $gx\in\mathfrak{p}_{m+n}$, ce qui prouve que $\mathfrak{p}$ est un
idéal de $S$. Pour établir que cet idéal est premier, autrement dit que
l'anneau $S/\mathfrak{p}$, gradué par les sous-groupes
$S_n/\mathfrak{p}_n$, est un anneau intègre, il suffit (en considérant les
composants de plus haut degré de deux éléments de $S/\mathfrak{p}$) de
prouver que si $x\in S_m$, $y\in S_n$ sont tels que
$x\notin\mathfrak{p}_m$, $y\notin\mathfrak{p}_n$, alors
$xy\notin\mathfrak{p}_{m+n}$. Sinon, pour $r$ assez grand, on aurait
$f^{2r}xy\in\mathfrak{p}_{m+n+2rd}$~; mais on a
$f^ry\notin\mathfrak{p}_{n+rd}$ pour tout $r>0$~; il résulte alors de
$2^\circ$ que, sauf pour un nombre fini de valeurs de $r$, on a
$f^rx\in\mathfrak{p}_{m+rd}$, et on en conclurait que
$x\in\mathfrak{p}_m$ contrairement à l'hypothèse.
\end{proof}

\begin{env}[2.1.10]
\label{II.2.1.10-fr}
Nous dirons qu'une partie $\mathfrak{J}$ de $S_+$ est un
\emph{idéal de $S_+$} si c'est un idéal de $S$, et que $\mathfrak{J}$ est un
\emph{idéal premier gradué de $S_+$} si elle est l'intersection de $S_+$ et
d'un idéal premier gradué de $S$ \emph{ne contenant pas $S_+$} (cet idéal
premier est d'ailleurs unique d'après
(\hyperref[II.2.1.9-fr]{2.1.9})). Si $\mathfrak{J}$ est un idéal de $S_+$,
la \emph{racine de $\mathfrak{J}$ dans $S_+$} est l'ensemble des éléments de
$S_+$ dont une puissance appartient à $\mathfrak{J}$, autrement dit c'est
l'ensemble $r_+(\mathfrak{J})=r(\mathfrak{J})\cap S_+$~; en particulier, la
racine de $0$ dans $S_+$ est encore appelée le \emph{nilradical} de $S_+$ et
notée $\mathfrak{N}_+$~: c'est donc l'ensemble des éléments nilpotents de
$S_+$. Si $\mathfrak{J}$ est un idéal \emph{gradué} de $S_+$, sa racine
$r_+(\mathfrak{J})$ est un idéal gradué~: en passant à l'anneau quotient
$S/\mathfrak{J}$, on peut se ramener au cas $\mathfrak{J}=0$, et tout
revient à voir que si $x=x_h+x_{h+1}+\cdots+x_k$ est nilpotent, il en est de
même des $x_i\in S_i$ ($1\leq h\leq i\leq k$)~; on peut supposer
$x_k\neq 0$ et le composant de plus grand degré de $x^n$ est alors $x_k^n$,
donc $x_k$ est nilpotent, et on raisonne alors par récurrence sur $k$. On dit
que l'anneau gradué $S$ est \emph{essentiellement réduit} si
$\mathfrak{N}_+=0$, autrement dit si $S_+$ ne contient pas d'éléments
nilpotents $\neq 0$.
\end{env}

\begin{env}[2.1.11]
\label{II.2.1.11-fr}
On notera que si, dans l'anneau gradué $S$, un élément $x$ est diviseur de
$0$, il en est de même de son composant de plus haut degré. On dit qu'un
anneau $S$ est \emph{essentiellement intègre} si l'anneau $S_+$
(\emph{sans élément unité}) ne contient pas de diviseur de $0$ et est
$\neq 0$~; il suffit pour cela qu'un élément homogène $\neq 0$ dans $S_+$ ne
soit pas diviseur de $0$ dans cet anneau. Il est clair que si $\mathfrak{p}$
est un idéal premier gradué de $S_+$, $S/\mathfrak{p}$ est essentiellement
intègre.

Soit $S$ un anneau gradué essentiellement intègre, et soit $x_0\in S_0$~;
s'il existe alors \emph{un} élément homogène $f\neq 0$ de $S_+$ tel que
$x_0f=0$, on a $x_0S_+=0$, car on a $(x_0g)f=(x_0f)g=0$ pour tout
$g\in S_+$, et l'hypothèse entraîne donc $x_0g=0$. Pour que $S$ soit
intègre, il faut et il suffit donc que $S_0$ soit intègre et que l'annulateur
de $S_+$ dans $S_0$ soit réduit à $0$.
\end{env}

\subsection{Anneaux de fractions d'un anneau gradué.}
\label{subsection:II.2.2-fr}

\begin{env}[2.2.1]
\label{II.2.2.1-fr}
Soient $S$ un anneau gradué, à degrés positifs, $f$ un élément
\emph{homogène} de $S$, de degré $d>0$~; alors l'anneau de fractions
$S'=S_f$ est gradué, en prenant pour $S'_n$ l'ensemble des $x/f^k$, où
$x\in S_{n+kd}$ avec $k\geq 0$ (on observera ici que $n$ peut prendre des
valeurs négatives arbitraires)~; nous désignerons le sous-anneau
$S'_0=(S_f)_0$ de $S'$ formé des éléments \emph{de degré $0$} par la
notation $S_{(f)}$.

Si $f\in S_d$, les monômes $(f/1)^h$ dans $S_f$ ($h$ entier positif ou
négatif) forment un \emph{système libre} sur l'anneau $S_{(f)}$, et
l'ensemble de leurs combinaisons linéaires n'est autre que
